\documentclass[10pt]{article}
\usepackage[margin=1in]{geometry}
\usepackage[T1]{fontenc}
\usepackage{lmodern,amsmath,amssymb,amsthm,microtype}
\usepackage[colorlinks=true,allcolors=blue]{hyperref}
\newtheorem{theorem}{Theorem}
\newtheorem{lemma}{Lemma}
\newtheorem{remark}{Remark}
\newcommand{\im}{\operatorname{im}}
\newcommand{\Span}{\operatorname{span}}
\newcommand{\dual}{\mathrm D}
\title{A counterexample to Takao's self-duality question}
\author{Alec Kriebel\\Independent researcher\\
\small ORCID: \href{https://orcid.org/0009-0001-9320-500X}{0009-0001-9320-500X}}
\date{4 October 2026}
\begin{document}
\maketitle
\begin{abstract}
Takao asked whether a finite flat commutative group scheme over a Witt
ring, killed by $p$ and having quasi-supersingular special fiber, must be
Cartier self-dual in ranks $p^{2n}$ with $n\geq3$. For every $p>3$ and
$n\geq3$, we give a counterexample over $W(\overline{\mathbb F}_p)$.
A six-dimensional Dieudonn\'e module has an explicit filtration by
supersingular elliptic $p$-torsion modules, an explicit finite Honda
complement, and an image-intersection invariant distinguishing it from
its dual. Both the special fiber and its finite flat Witt-ring lift
are non-self-dual. The construction is an application of classical
Dieudonn\'e and finite Honda theory.
\end{abstract}

\section{The question and result}
Let $p>3$ be prime. A finite commutative group scheme $H$ over a perfect
field of characteristic $p$, killed by $p$, is \emph{quasi-supersingular}
if it has a subgroup filtration whose successive quotients are
$E_i[p]$ for supersingular elliptic curves $E_i$ over that field.
This is Takao's Definition~1(3)~\cite{Takao}. His Proposition~2 states
that, for a finite flat commutative $W(k)$-group scheme $\mathcal G$
killed by $p$, of rank $p^{2n}$, a quasi-supersingular special fiber
forces both the special fiber and $\mathcal G$ to be self-dual when
$n\leq2$. The unnumbered question immediately following that proposition
asks whether the assertion remains true for $n\geq3$
\cite[p.~2479]{Takao}.

\begin{theorem}\label{thm:main}
For every prime $p>3$ and integer $n\geq3$, put
$k=\overline{\mathbb F}_p$ and $W=W(k)$.
There is a finite flat commutative group scheme $\mathcal G_n/W$,
killed by $p$ and of rank $p^{2n}$, whose special fiber $H_n$ is
quasi-supersingular, such that
\[
 H_n\not\simeq H_n^{\dual},\qquad
 \mathcal G_n\not\simeq\mathcal G_n^{\dual}.
\]
Here ${}^{\dual}$ denotes Cartier duality.
\end{theorem}

Thus the answer to the stated higher-rank question is negative, with
the first possible counterexample at $n=3$, using Takao's affirmative
result for $n\leq2$. The hypothesis concerns the special fiber: we do
not require its filtration to lift to $W$.

The module used below belongs to the classical cyclic-word framework.
In the convention of Pries--Ulmer it is $M(vvfvff)$
\cite[\S\S3.2.1,4.1]{PriesUlmer}; in the forward-arrow order of our
contravariant presentation its string is $ffvfvv$.
The module construction and its duality classification are prior
mathematics. Oort's cyclic integral construction also supplies a
balanced supersingular completion~\cite[\S\S2.5--2.6]{Oort}.
We give direct calculations of the filtration, duality obstruction,
and finite Honda complement to answer the specific question without
relying on the cyclic-word classification. No first-discovery or
first-priority claim is made. The supporting package records the
bounded literature audit and its source-access and version limits.

\section{The six-dimensional construction}
We use contravariant Dieudonn\'e theory over $k$ as in
\cite[\S2]{Hoshi}. Write $\sigma$ for Frobenius on $k$ and also for
its lift to $W$. On a module killed by $p$, $F$ is
$\sigma$-semilinear, $V$ is $\sigma^{-1}$-semilinear, and $FV=VF=0$.
Let $M$ have basis $e_0,\ldots,e_5$, with operators specified by
\[
\begin{array}{c|rrrrrr}
 &e_0&e_1&e_2&e_3&e_4&e_5\\\hline
F&e_1&e_2&0&e_4&0&0\\
V&e_5&0&0&e_2&0&e_4
\end{array}
\]
and extended with the indicated semilinearity. All displayed scalars
are in $\mathbb F_p$. Direct computation gives
\begin{align*}
 FV=VF&=0,& F^3=V^3&=0,\\
 \im F=\ker V&=\Span_k(e_1,e_2,e_4),&
 \im V=\ker F&=\Span_k(e_2,e_4,e_5).
\end{align*}
The finite Dieudonn\'e anti-equivalence produces a commutative
$k$-group scheme $H$ killed by $p$, of rank $p^6$
\cite[Proposition~2.5]{Hoshi}.

\subsection{Quasi-supersingularity}
Let $I$ denote the two-dimensional module on $x,y$ with
\[
 Fx=Vx=y,\qquad Fy=Vy=0.
\]
Over algebraically closed $k$, $I$ is the module of the $p$-torsion
of a supersingular elliptic curve. This is the rank-two case of
\cite[Lemma~4.9]{Hoshi}: the module is deformable and
$\im F=\im V$.

Set
\[
 u=e_1+e_3+e_5,\quad v=e_2+e_4,\quad
 w=2e_1+e_3,\quad z=e_2,\quad a=e_0,\quad b=e_1.
\]
These form a basis. Indeed,
\[
 e_0=a,\ e_1=b,\ e_2=z,\ e_3=w-2b,\
 e_4=v-z,\ e_5=u-w+b.
\]
In this basis the nonzero images are exactly
\begin{align*}
 Fu&=v,& Fw&=v+z,& Fa&=b,& Fb&=z,\\
 Vu&=v,& Vw&=z,& Va&=u-w+b.
\end{align*}
Consequently
\[
 0=M_0\subset M_1=\Span_k(u,v)
 \subset M_2=\Span_k(u,v,w,z)\subset M_3=M
\]
is an $F,V$-stable flag. The three quotient modules are $I$, using,
respectively, the pairs $(u,v)$, $(w,z)$, and $(a,b)$.
Exact contravariance gives subgroups $H_i\subset H$ with modules
$M/M_{3-i}$, for $0\leq i\leq3$.
Their successive quotients correspond in reverse order to the three
module factors. Each is supersingular elliptic $p$-torsion; hence
$H$ is quasi-supersingular.

\subsection{Failure of Cartier self-duality}
For a finite Dieudonn\'e module $T$ over perfect $k$, define
\[
 \delta(T)=\dim_k(\im F_T^2\cap\im V_T^2).
\]
The images are $k$-subspaces because the field is perfect.
An isomorphism commuting with $F,V$ carries both images to the
corresponding images, so $\delta$ is an isomorphism invariant.
For $M$,
\[
 \im F^2=k e_2,\qquad \im V^2=k e_4,\qquad \delta(M)=0.
\]
On $M^{\dual}=\operatorname{Hom}_k(M,k)$ the dual operators are
\[
 (F^{\dual}\phi)(t)=\sigma\bigl(\phi(Vt)\bigr),\qquad
 (V^{\dual}\phi)(t)=\sigma^{-1}\bigl(\phi(Ft)\bigr)
 \quad\text{\cite[Definition~2.3]{Hoshi}}.
\]
Since our basis matrices have coefficients in $\mathbb F_p$, their
dual matrices are $V^{\mathsf t}$ and $F^{\mathsf t}$ with the
appropriate semilinearity. Thus $(F^{\dual})^2(e_4^*)=e_0^*$ and
$(V^{\dual})^2(e_2^*)=e_0^*$, and all other basis images of these
second iterates vanish. It follows that
\[
 \im(F^{\dual})^2=\im(V^{\dual})^2=k e_0^*,\qquad
 \delta(M^{\dual})=1.
\]
Hence $M\not\simeq M^{\dual}$, and compatibility of Dieudonn\'e
theory with Cartier duality gives $H\not\simeq H^{\dual}$.
The transpose description here uses the prime-field coefficients;
the displayed semilinear dual formulas are the general convention.

\section{An actual finite flat lift and all higher ranks}
View $M$ as a $W$-module through $W\to k$ and put
\[
 L=\Span_k(e_0,e_3,e_5)\subset M.
\]
We verify the finite Honda hypotheses explicitly. The $W$-length of
$M$ is six and $pM=0$. The semilinear operators satisfy $FV=VF=p$
on $M$. The restriction of $V$ to $L$ sends its displayed basis
to $e_5,e_2,e_4$ and is injective, since $k$ is perfect. Finally,
$L$ complements $\im F=\Span_k(e_1,e_2,e_4)$, so the natural map
\[
 L/pL\longrightarrow M/\im F
\]
is an isomorphism. These are exactly the finite Honda conditions
of Fontaine--Laffaille~\cite[\S9.4]{FontaineLaffaille}, in the
$p$-torsion formulation of Hoshi~\cite[Remark~3.5.1]{Hoshi}.
The exactness identities above also verify Hoshi's deformability
condition.

Since $p\ne2$, Hoshi's anti-equivalence
\cite[Proposition~3.11(2)]{Hoshi} produces a finite flat commutative
$W$-group scheme $\mathcal G$, killed by $p$, from this finite Honda
system. Its special fiber has module $M$, by
\cite[Proposition~3.11(5)]{Hoshi}, and is therefore $H$.
Its rank is $p^6$: finite flat rank is constant over the local base
and equals the special-fiber rank. Cartier duality commutes with
base change, so an isomorphism $\mathcal G\simeq\mathcal G^{\dual}$
would imply $H\simeq H^{\dual}$. Thus $\mathcal G$ also fails
self-duality. This is a finite Honda realization over $W$; an
integral Dieudonn\'e lattice for a group over $k$ would not, on
its own, establish this mixed-characteristic assertion.

For $n\geq3$, equip $I$ with the finite Honda subspace $kx$ and set
\[
 M_n=M\oplus I^{\oplus(n-3)},\qquad
 L_n=L\oplus(kx)^{\oplus(n-3)}.
\]
The same Honda conditions hold on each block. They give
$\mathcal G_n/W$ of rank $p^{2n}$. Its special fiber is
$H\times E[p]^{n-3}$ and has a quasi-supersingular filtration by
extending the filtration already constructed with the product
elliptic factors. Since $F^2=V^2=0$ on $I$ and on its dual,
\[
 \delta(M_n)=0,\qquad\delta(M_n^{\dual})=1.
\]
Therefore $H_n$ and, by base change, $\mathcal G_n$ are not
Cartier self-dual. This proves Theorem~\ref{thm:main}.

\paragraph{Scope and verification.}
These examples over $\overline{\mathbb F}_p$ disprove the universal
assertion for perfect residue fields; no descent to every perfect
field is asserted. The result concerns general finite group schemes
and does not produce a principally polarized Jacobian or decide
the surrounding Coleman conjecture. The proof is symbolic for
every allowed $p,n$; the companion standard-library verification
programs check the basis, operators, invariant, Honda conditions,
direct sums, and nontrivial Frobenius twists over finite fields.
Those computations corroborate the proof and do not replace the
cited classification theorems or certify historical priority.

\paragraph{AI use and review status.}
AI tools were used extensively in solving, verification, literature
research, writing, computation, and adversarial review. This note
is an unrefereed preprint. The automated reviews are not independent
external human peer review; no such human review is claimed.

\begin{thebibliography}{9}
\bibitem{FontaineLaffaille}
J.-M. Fontaine and G. Laffaille,
\emph{Construction de repr\'esentations $p$-adiques},
Ann. Sci. \'Ecole Norm. Sup. (4) \textbf{15} (1982), no.~4,
547--608,
\href{https://www.numdam.org/item/ASENS_1982_4_15_4_547_0/}{Numdam original}.
\bibitem{Hoshi}
Y. Hoshi, \emph{Pseudo-rigid $p$-Torsion Finite Flat Commutative
Group Schemes}, March 2021 revised author manuscript,
\url{https://www.kurims.kyoto-u.ac.jp/~yuichiro/rims1911revised.pdf}.
The numbered statements cited are those of this manuscript.
Published in J. Number Theory \textbf{229} (2021), 261--276,
\href{https://doi.org/10.1016/j.jnt.2021.04.026}{doi:10.1016/j.jnt.2021.04.026}.
\bibitem{Oort}
F. Oort, \emph{Simple $p$-kernels of $p$-divisible groups},
June 2004/revised January 2005 author manuscript,
\url{https://webspace.science.uu.nl/~oort0109/A-SimpleFinGrSch2.ps}.
Published in Adv. Math. \textbf{198} (2005), 275--310,
\href{https://doi.org/10.1016/j.aim.2005.03.015}{doi:10.1016/j.aim.2005.03.015}.
\bibitem{PriesUlmer}
R. Pries and D. Ulmer, \emph{On $\mathrm{BT}_1$ group schemes and
Fermat curves}, New York J. Math. \textbf{27} (2021), 705--739,
\url{https://nyjm.albany.edu/j/2021/27-27v.pdf}.
\bibitem{Takao}
N. Takao, \emph{Quasi-supersingular finite flat commutative group
schemes and the Coleman conjecture on torsion points on curves},
in \emph{Arithmetic Homotopy and Galois Theory},
Oberwolfach Reports \textbf{20} (2023), no.~3,
Report~42/2023, 2476--2480,
\href{https://doi.org/10.4171/OWR/2023/42}{doi:10.4171/OWR/2023/42}.
\end{thebibliography}
\end{document}
